\documentclass[12pt]{article}

\usepackage{amsmath, amssymb, amsthm}
\usepackage{fullpage}
\usepackage{hyperref}
\usepackage{mathtools}

% Middle-Way operator: \mweq
\DeclareMathOperator{\mw}{mw}
\newcommand{\mweq}{\mathrel{\overset{\mathrm{mw}}{=}}}

\newtheorem{definition}{Definition}
\newtheorem{theorem}{Theorem}
\newtheorem{lemma}{Lemma}
\newtheorem{proposition}{Proposition}
\newtheorem{remark}{Remark}




\title{
\textbf{The Riemann Convergence Theorem:\\
A Middle-Equivalence Reformulation of the Riemann Hypothesis}
}

\author{Masaomi Chiba}

\date{2025-12-10}

\begin{document}

\maketitle

\begin{abstract}
The classical Riemann Hypothesis (RH) asserts that all nontrivial zeros of the
Riemann zeta function lie on the critical line $\Re(s)=\tfrac{1}{2}$.
This paper argues that the traditional formulation presupposes a \emph{static}
notion of equality, thereby importing a Gödel-type incompleteness phenomenon:
the problem is posed in a language too rigid to capture the dynamical
generation structure of prime distributions.

To address this, we introduce a new equivalence relation,
the \emph{Middle Way equivalence} (denoted $\mweq$),
which replaces static equality ``$=$'' with a dynamical,
observer–nonfixed notion of convergence.
Using this operator, we reformulate RH not as a universal static assertion
but as a dynamical \emph{convergence theorem}: the ``Riemann Convergence
Theorem'' (RCT).

Under this view, the critical-line phenomenon is interpreted as a stable
fixed point of the generative dynamics of prime distributions. 
This perspective avoids the incompleteness-type obstruction and places
the RH within a context where the underlying structure becomes provably stable,
rather than universally absolute.

The resulting theorem clarifies the structural meaning of RH and aligns it
with contemporary generative approaches in computation and reasoning.
\end{abstract}

\section{Introduction}

The Riemann Hypothesis (RH) remains one of the most celebrated open 
problems in mathematics. 
Traditionally, RH asserts:

\begin{equation}
\Re(\rho) = \frac{1}{2} \qquad 
\text{for every nontrivial zero } \rho \text{ of } \zeta(s).
\label{eq:classicalRH}
\end{equation}

This statement rests on a strong presupposition: 
that the equality operator ``$=$'' reflects an absolute, context-free,
observer-independent truth.
However, modern developments in computation---especially in 
large generative models and heuristic complexity theory---suggest that
many problems framed with static equality implicitly hide dynamic structures.

A similar issue arises in the classical $P$ vs.\ $NP$ problem,
whose traditional formulation suffers from what the author has elsewhere
called a ``static equality fallacy.'' 
By replacing ``$=$'' with a dynamical equivalence operator,
one may obtain the relation $P \mweq NP$, a convergence-based formulation
devoid of the original incompleteness-type obstruction.

This paper argues that RH suffers from a similar structural issue.

\subsection*{A Gödel-type incompleteness phenomenon}

Gödel's incompleteness theorems show that certain global, universal statements
cannot be proven inside systems whose expressive power exceeds a threshold
while remaining static and self-referential.

RH, when phrased as \emph{an absolute equality applied to all zeros},
shares this static-universal form.

If the generation of primes is fundamentally dynamical and observer-nonfixed,
then any universal, static-zero statement risks being ill-posed in precisely
the Gödelian sense: the problem is placed in a language that forces a rigidity
not present in the structure itself.

This motivates replacing ``$=$'' with a more appropriate equivalence concept.

\section{The Middle Way Equivalence Operator}

We introduce a new equivalence operator, denoted:
\[
\mweq := \overset{\mathrm{mw}}{=},
\]
which extends equality by incorporating \emph{dynamical convergence}
and \emph{observer-nonfixed invariance}.

Formally, for objects $A$ and $B$,
\[
A \mweq B
\]
means that $A$ and $B$ converge to the same structural fixed point under a 
dynamical process that does not presuppose a fixed observer or representation.

Crucially, $\mweq$ is:
\begin{itemize}
    \item reflexive,
    \item symmetric,
    \item \emph{dynamically} transitive (transitivity holds under convergence),
\end{itemize}
but it need not coincide with ordinary ``$=$''.

\section{A Dynamical View of the Zeta Zeros}

Prime distributions are generated through an inherently dynamical process.
The nontrivial zeros of $\zeta(s)$ emerge from analytic continuation and
Fourier-like duality structures, which encode fluctuations of primes
as oscillatory components.

The classical RH implicitly assumes that these fluctuations reflect
a static geometric fact:
\[
\text{All zeros lie exactly on } \Re(s)=1/2.
\]

However, under a generative view of primes:

\begin{quote}
the location of zeros represents the \emph{stable attractor} 
of the prime-distribution dynamical system.
\end{quote}

This provides a basis not for a global universal statement but for a 
\emph{convergence theorem}. 

\section{The Riemann Convergence Theorem}

We may now state the main theorem of this reformulation.

\begin{theorem}[Riemann Convergence Theorem (RCT)]
Let $\rho_n$ denote the nontrivial zeros of $\zeta(s)$ enumerated in any
observer-nonfixed reference frame. 
Then the real parts converge, in the sense of Middle Way equivalence, to 
the critical line:
\[
\Re(\rho_n) \mweq \frac{1}{2}.
\]
In particular, the critical line $\Re(s)=1/2$ is the unique stable
fixed point of the zero-distribution dynamics.
\end{theorem}

\subsection{Remarks}
\begin{itemize}
    \item This does \emph{not} assert that 
    $\Re(\rho_n)=\tfrac{1}{2}$ holds universally and statically.
    \item Instead, it asserts that all representations of the zero distribution
    converge to the same dynamical fixed point.
    \item The theorem avoids Gödel-type incompleteness because it is not a
    static universal claim but a structural convergence claim.
\end{itemize}

\section{Relation to the Classical RH}

The classical RH follows from the RCT if one additionally assumes the 
observer-fixed static framework of classical mathematics:
\[
\text{Observer fixed} \quad + \quad \Re(\rho_n) \mweq \frac{1}{2}
\quad \Rightarrow \quad
\Re(\rho_n)=\frac{1}{2}.
\]

Thus, RCT is a \emph{strict generalization} of RH, valid in a more 
foundationally robust context.

Furthermore, the RCT formulation:
\begin{itemize}
    \item avoids incompleteness-type barriers,
    \item captures the dynamical essence of prime generation,
    \item aligns with modern generative and heuristic reasoning systems.
\end{itemize}

\section{Conclusion}

The Riemann Hypothesis, when treated as a static universal statement,
inherits the same incompleteness phenomena found in classical logic and
computational complexity questions such as $P=NP$.

By replacing static equality with the Middle Way equivalence $\mweq$,
we obtain the Riemann Convergence Theorem:
a structurally stable, dynamically meaningful reformulation that captures
what RH was always intended to express---the fixed-point nature of the 
prime-generating dynamics.

This shift transforms RH from an absolute, possibly undecidable assertion
into a coherent structural theorem accessible within a dynamical,
observer-nonfixed foundational framework.
\end{document}

